\documentclass[./../main.tex]{subfiles} % 继承主文件设置
\begin{document}
\section{chapter10: 随机算法}
\subsection{导论}


\begin{enumerate}
    \item 与确定性算法相比，随机算法所需运行时间或空间通常较小；
    \item 随机算法易于理解和实现。
\end{enumerate}

\begin{enumerate}
    \item \textbf{Las Vegas 算法}\footnote{Las Vegas 和 Monte Carlo 不是一种算法的名称，而是对一类随机算法的特性的概括。不太可能为\textbf{NP完全问题设计出多项式时间的随机算法}}：总是给出正确的解，或者无解

    通常需要分析\textbf{算法的平均运行时间}，运行时间本身是一个随机变量。

   \textbf{期望运行时间是输入规模的多项式}且\textbf{总能给出正确答案}的随机算法称为\textbf{有效}的拉斯维加斯算法，也称为 \textbf{ZPP 类算法}(zero-error probabilistic polynomial time)算法，即零错误概率多项式时间随机算法。
   
    \item \textbf{Monte Carlo 算法}：总是给出解，偶尔可能是错误解，可以通过多次运行原算法，且满足每次运行时随机选择都相互独立，使产生错误解的概率减到任意小。

    通常需要分析\textbf{算法的出错概率}，总是在\textbf{多项式时间内运行}且出错概率\textbf{不超过 1/3}的随机算法称为\textbf{有效的*}蒙塔卡洛型随机算法，也称为\textbf{BPP 类算法}(bounded-error probabilistic polynomial time)，即错误概率有限的多项式时间随机算法。

    错误类型：
    \begin{itemize}
        \item \textbf{弃真型错误}：在求解一个判定问题时把本应该接受的输入误判为拒绝。
        \item \textbf{取伪型错误}：就是把本应该拒绝的输入误判为接受。
    \end{itemize}

    根据是否同时出现两类错误，蒙特卡洛型算法又分为：
    \begin{enumerate}
        \item \textbf{单侧错误的蒙特卡洛型算法}：在所有输入上，要么犯弃真的错误，要么犯取伪的错误，不同时出现两种不同的错误。
        \begin{itemize}
            \item 有效的\textbf{弃真型}单侧错误随机算法称为\textbf{RP 类算法}(Randomized Polynomial-time)
            \item 有效的\textbf{取伪型}单侧错误随机算法称为\textbf{coRP 类算法}(Complementary Randomized Polynomial-time, RP 的补类)
        \end{itemize}
        \item \textbf{双侧错误}的蒙特卡洛型算法：在所有输入上可以同时出现两种不同的错误。
    \end{enumerate}
\end{enumerate}

对于随机算法 A, 其在大小为 n 的一个实例 I 上的运行时间可能因不同次执行而异。因此，更自然的性能度量是 \textbf{A 在实例 I 上的期望运行时间}，也即算法 A 反复求解实例 I 所花费的平均时间。

\begin{itemize}
    \item 对于确定算法，求解期望运行时间时，概率分布定义在\textbf{输入}上
    \item 对于随机算法，求解期望运行时间时，概率分布定义在\textbf{算法}上
\end{itemize}

期望时间：对于大小为 n 的特定输入I, 算法的期望时间就是在随机数的不同序列上的时间的平均: 

\begin{equation*}
    ET(A,I)=\sum_{\text{random sequence R}}\text{probility}(R)*\text{time}(A,I,R)
\end{equation*}

最坏情况下的期望时间：

\begin{equation*}
   W\_ET(A)=\max_{\text{input I with size n}}(\text{random sequence R}\text{probility}(R)*\text{time}(A,I,R))
\end{equation*}

\subsection{随机算法举例}

这是一个很不错的讲解栏目：\url{https://www.luogu.com.cn/article/wegx023s}

\subsubsection{随机快速排序}

其实是在快速排序的基础上，增加了\textbf{随机选取}的步骤。
\begin{algorithm}[t]
    \SetKwData{A}{A}\SetKwData{Low}{low}\SetKwData{High}{high}\SetKwData{W}{w}
    \SetKwFunction{Quicksort}{quicksort}\SetKwFunction{Split}{SPLIT}
    \SetKwInOut{Input}{Input}
    \SetKwInOut{Output}{Output}

    \Input{An array $A[1\ldots n]$ of $n$ elements}
    \Output{The elements in $A$ sorted in nondecreasing order}
    \BlankLine

    \Quicksort(\A,\Low,\High)\\
    \If {\Low < \High} {
        \Split{\A[\Low...\High], \W}\\
        \tcp{\W is the new position of \A[\Low]}
        \Quicksort(\A,\Low, \W-1)\\
        \Quicksort(\A, \W+1, \High)
    }
    \caption{Quicksort}
\end{algorithm}

\begin{algorithm}[t]
    \SetKwData{A}{A}\SetKwData{low}{low}\SetKwData{high}{high}\SetKw{w}{w}\SetKw{v}{v}
    \SetKwFunction{rquicksort}{rquicksort}
    \SetKwFunction{SPLIT}{SPLIT}
    \SetKwFunction{random}{random}
    \SetKwInOut{Input}{Input}\SetKwInOut{Output}{Output}

    \Input{An array $A$ of $n$ elements}
    \Output{The elements in $A$ sorted in nondecreasing order}
    \BlankLine

    \rquicksort(\A,\low,\high)

    \If {\low < \high} {
        \v $\leftarrow$ \random(\low, \high)
        \emph{interchange} \A[\low] \emph{and} \A[\v]\\
        \SPLIT{\A[\low...\high], \w, \v}\\
        \rquicksort(\A,\low, \w-1)\\
        \rquicksort(\A,\w+1, \high)\\
    }
    \caption{Randomized Quicksort}
\end{algorithm}

该算法在最坏情况的运行时间仍然是$\Theta(n^2)$ ，但这个最坏情况不是由于输入形式，而是由于随机数生成器每次非常不幸地选择了不好的划分元素，这种情况不大可能发生。算法的期望运行时间为$\Theta(n\log n)$, 是一个 ZPP 类算法

\subsubsection{随机选择}

\begin{algorithm}[!t]
    \SetKwData{high}{high}\SetKwData{low}{low}\SetKwData{mid}{mid}\SetKwData{pivot}{pivot}\SetKwData{q}{q}\SetKwData{p}{p}\SetKwData{n}{n}\SetKwData{k}{k}\SetKwData{A}{A}\SetKwData{mm}{mm}
    \SetKwFunction{Select}{select}
    \SetKwInOut{Input}{Input}\SetKwInOut{Output}{Output}

    \Input{An array $A$ of $n$ elements, a number $k$ such that $1\leq k\leq n$}
    \Output{The $k$th smallest element in $A$}
    \Select{\A, 1, n, \k}

    \BlankLine
    \p $\leftarrow$ \high - \low + 1\\
    \If{\p < 44}{sort \A and return \A[\k]}
    let \q = $\lfloor p/5 \rfloor$.Divide \A into \q groups of 5 elements each. If 5 doesn't divide \p, then discard the remaining elements.\\
    \tcp{Although some elements are discarded here, the subsequent grouping still takes them into account, so it will not affect the answer.}
    Sort each of the \q groups individually and extract its median element. Let the set of medians be $M$ \\ 
    \mm $\leftarrow$ \Select($M$, 1, \q, $\lfloor$\q/2$\rfloor$) \tcp{median of medians}
    Partition \A into three arrays:\\
    $A_1\leftarrow \{A|A<mm\}$ ,$A_2\leftarrow\{A|A=mm\}$,$A_3 \leftarrow\{A|A>mm\}$\\
    \eIf{$|A_1|\ge k $}{return \Select($A_1$, 1, $|A_1|$, \k)}{
        \eIf{$k\le |A_1|+|A_2|$}{return mm}{return \Select($A_3$, 1, $|A_3|$, \k$-|A_1|-|A_2|$)}
    }
    
    \caption{SELECT algorithm}
\end{algorithm}


\begin{algorithm}[t]
    \SetKwFunction{rselect}{rselect}\SetKwFunction{random}{random}

    \KwIn{An array $A$ of $n$ elements, a number $k$ such that $1\leq k\leq n$}
    \KwOut{The $k$th smallest element in $A$}
    \rselect{$A, 1, n, k$}

    $v\leftarrow $ \random{$low, high$}\\
    $x\leftarrow A[v]$\\
    Partition $A[low...high]$ into three arrays: $A_1\leftarrow\{A|A<x\}$, $A_2\leftarrow\{A|A=x\}$, $A_3\leftarrow\{A|A>x\}$

    \eIf{$|A_1|\ge k $}{return \rselect{$A_1$, 1, $|A_1|$, $k$}}{
        \eIf{$k\le |A_1|+|A_2|$}{return x}{return \rselect{$A_3$, 1, $|A_3|$, $k-|A_1|-|A_2|$}}
    }
    \caption{Randomized SELECT}
\end{algorithm}

在最坏情况下，随机数生成器在第$j$次迭代中选择第$j$小的元素，因此只丢弃一个元素。由于将$A$划分为$A_1$、$A_2$和$A_3$所需的比较次数正好为$n$，因此算法在最坏情况下进行的元素比较总数为
$$
\sum_{i=1}^n i=\frac{n(n+1)}{2}=\Theta(n^2)\\
$$
设$C(n)$为对输入数组有$n$个元素时进行元素比较运算的\textbf{期望次数}。由于随机选择的v可以等概率假设为$1,2,...,n$中的任何一个，因此根据$v<k$或$v>k$，有两种情况需要考虑。
\begin{itemize}
    \item $v<k$: 剩余元素的数量为$n-v$($v$右侧的元素数量为$n-v$)
    \item $v>k$: 剩余元素的数量为$v-1$($v$左侧的元素数量为$v-1$)
\end{itemize}

因而有：
$$
C(n)=n+\frac{1}{n}\sum_{v=1}^n\left\{
    \begin{array}{lc}
        C(v-1) & \text{if }v<k\\ 
        C(n-v) & \text{if }v>k\\
    \end{array}
\right.
$$

由数学归纳法，假设$C(m)\le4m$对所有$m<n$成立，下证$C(n)\le4n$:
\begin{align*}
    C(n)&\le n + \frac{1}{n}\sum_{v=1}^n4\cdot\max(n-v, v-1)\\
    &\le n + \frac{8}{n}\cdot \sum_{v=1}^{\frac{n}{2}}(n-v)\\
    &\le n + \frac{8}{n}\cdot \frac{3n^2}{8}\\
    &=  4n\\
\end{align*}
又由于每个元素至少被检查一次，因此$C(n)\ge n$，\textbf{因此算法的期望运行时间是}$\Theta(n)$，是 ZPP 类算法。

\subsubsection{多选问题}
设有$n$个元素的数组$A$, 并且设$K[1..r](1\le r\le n)$是一个由1到$n$之间的$r$个正数构成的有序数组，即一个排序数组。

多选问题就是对于所有$i$，$1\le i\le r$，选择A中的第$K[i]$小的元素。假定$A[1...n]$中所有元素各不相同。如果$r=1$，该问题即为选择问题。如果$r=n$，该问题即为排列问题。

设$A[1..n]=[a_1, a_2, \dots, a_n]; K[1..r]=[k_1, k_2, \dots, k_r]$，设中间排位为$k=k_{\frac{r}{2}}$,可使用SELECT算法发现第$k$小的元素$a$.
接着根据$a$把$A$划分成两个子序列$A_1$和$A_2$，其中$A_1$中元素个数大于等于$a$，$A_2$中元素个数小于等于$a$。令$K_1=[k_1,k_2,...k_{\frac{r}{2}}], K_2=[k_{\frac{r}{2}+1},k_{\frac{r}{2} + 2},...,k_r]$。
递归调用，一个关于$A_1$和$K_1$的递归调用，一个关于$A_2$和$K_2$的递归调用。

\begin{algorithm}
    \SetAlgoLined
    \KwIn{$n$元序列$A$，具有$r$个排位的有序序列$K$}
    \KwOut{$A$中第$k_i$小的元素}
    \BlankLine
    multiselect($A$, $K$)\\
    $r\leftarrow |K|$\\
    \If{$r>0$}{
        $k=k_{\frac{r}{2}}$\\
        $a\leftarrow$SELECT($A$, $k$)\\
        $Ans[\frac{r}{2}]\leftarrow   a$\\
        $A_1=[a_i|a_i<a], A_2=[a_i|a_i>a]$\\
        $K_1\leftarrow[k_1,k_2,...,k_{\frac{r}{2}}],K_2\leftarrow[k_{\frac{r}{2}+1}-k,...,k_r-k]$\\
        \tcp{特别注意这里需要减去$k$}
        multiselect($A_1$, $K_1$)\\
        multiselect($A_2$, $K_2$)\\
        \Return $Ans$
    }\caption{Multiselect}
\end{algorithm}
Multiselect算法的时间复杂度是$\Theta(n\log r)$.这是因为调用SELECT算法的复杂度是$\Theta(n)$, 所以有递推式
$$
T(n, r)= \Theta(n) + T(n_1, \frac{r}{2}) + T(n_2, \frac{r}{2}), n_1+n_2\le n
$$

递归深度是$\log r$，所以时间复杂度是$\Theta(n\log r)$.

将多选问题要寻找的元素称为靶标。
随机一致地挑选一个元素$a$（而不是通过SELECT精心选择），将数组$A$关于$a$分割成两部分：一部分元素小于$a$, 另一部分元素大于$a$
​处理逻辑：​​

如果两部分均含有靶标，则递归调用算法继续求解;
否则，丢弃不含有靶标的部分，对含有靶标部分继续求解

该算法以概率$1-O(\frac{1}{n})$的时间复杂度为$O(n\log r)$.
$r$个靶标的分配可以看作$r$次独立伯努利试验，其分布近似​二项分布$​B(r,1/2)$。
根据\emph{Hoeffding不等式}，靶标分配的偏离概率满足：
$$
P\left(\left|\text{左侧靶标数} - \frac{r}{2}\right| \geq \epsilon r\right) \leq 2\exp\left(-2\epsilon^2 r\right)
$$
当$r \geq c\log n$时，此概率上界为$O(1/n)$。

\subsubsection{随机取样}
首先将所有$n$个元素标记为未选中。
接下来，重复以下步骤，直到选择了$m$个元素：
\begin{enumerate}
    \item 生成一个介于$1$和$n$之间的随机数r
    \item 如果$r$被标记为未选中，则将其标记为选中并将其添加到样本中
\end{enumerate}

算法的期望运行时间$T(n)=\Theta(n)$
如果$m>\frac{n}{2}$，那么应该考虑反面，改为随机选择$n-m$个整数丢弃，将剩下的整数添加到样本中。

设第$k$次选择时已选中$k$个元素，则：
\begin{itemize}
    \item 成功选中新元素的概率 $p_k = \frac{n-k}{n}$
    \item 失败（重复选中）的概率 $q_k = 1-p_k = \frac{k}{n}$
\end{itemize}

根据几何分布，第$k+1$个新元素的期望尝试次数为：
$$
E_k = \frac{1}{p_k} = \frac{n}{n-k}
$$
每次尝试的时间成本为$\Theta(1)$（仅需检查布尔数组$S$）。

$$
\begin{aligned}
T(n) &= \sum_{k=0}^{m-1} E_k \cdot \Theta(1) \\
     &= \sum_{k=0}^{m-1} \frac{n}{n-k} \\
     &= n \left( \frac{1}{n} + \frac{1}{n-1} + \cdots + \frac{1}{n-m+1} \right) \\
     &= n \cdot (H_n - H_{n-m}) \quad \text{（调和数展开）}
\end{aligned}
$$

利用调和数近似 $H_n \approx \ln n + \gamma$：
$$
T(n) \approx n \left( \ln n - \ln(n-m) \right) = n \ln\left(\frac{n}{n-m}\right)
$$
\begin{itemize}
    \item 当$m = O(n)$时（如$m=\sqrt{n}$）：$T(n) \approx n \cdot \frac{m}{n} = \Theta(m)$($\ln x\approx x - 1$)
    \item 当$m = \Theta(n)$时（如$m=n/2$）：$T(n) = \Theta(n)$
\end{itemize}

\noindent 因此对于任意$m<n$，最坏情况下$T(n) = \Theta(n)$。


\subsubsection{素数测试}
易知除2，3以外的所有素数都是$6k\pm1$的形式，因此可以用$6k\pm 1\le \sqrt{n}$来对$n$是否是素数进行测试。

\begin{enumerate}
    \item \textbf{随机抽取一个数} $a$（$1 \leq a < n$）；
    \item \textbf{检查等式}：验证涉及$a$和$n$的特定数论等式（与所选测试方法对应）。\\
    若等式不成立，则：
    \begin{itemize}
        \item $n$ 是合数（Composite）
        \item $a$ 称为合数的\textbf{证据}（Witness）
        \item 测试终止
    \end{itemize}
    \item \textbf{重复迭代}：返回步骤1，直到达到预设置信度。
\end{enumerate}

现在我们先介绍幂运算：
\begin{algorithm}[h]
    \SetAlgoLined
    \KwIn{Positive integer $n$, positive integer $a$}
    \KwOut{Positive integer $a^n$}
    Let the binary digits of $n$ be $b_k, b_{k-1}, \dots, b_0$\\
    $c\leftarrow 1$\\
    \For{$j\leftarrow k$\textbf{downto} 0}{
        $c\leftarrow c\cdot 2$\\
        \If{$b_j$ is odd}{
            $c\leftarrow c\cdot a$
        }
    }
    \Return c
    \caption{EXP}
\end{algorithm}
如果每次乘法花费1个时间单元，算法的运行时间是$\Theta(\log n)$。但如果是任意大的整数，两数相乘的时间是$O(\log^2n)$，则算法的运行时间是$O(\log^3 n)$。

利用费马定理：如果$n$是素数，那么对于所有的$a\not\equiv0(mod~n)$，都有$a^{n-1}\equiv1(mod~n)$,如果通过$a$的测试，则成$n$是基2的伪素数，例如341.而对于小于100000的数中，只会有78个数计算出错，这是因为费马定理的逆并不正确。一类被称为\textbf{卡迈克尔数}的合数$n$,它对于所有相对于$n$互素的正整数$a$均满足费马定理，他们在小于1e8的范围内只有255个。

引理：如果合数$n$不是一个卡迈克尔数，那么算法（随机选择$a$并验证费马定理）将检测出$n$为合数的概率至少是$\frac{1}{2}$.

若$n$是合数但非卡迈克尔数，则存在至少一个$b\bot n$使得$b^{n-1}\not\equiv1(mod~n)$，则定义集合
$$
W=\{a\in\{1,2,...,n-1\}|a^{n-1}\not\equiv1(mod~n)\}
$$
而对于$W$在整数集$Z_n^*$上的补集$G$，易知$G$是$Z_n^*$的一个乘法子群，那么由拉格朗日定理知$|G| $整除$ |Z_n^*|=\phi(n)$，所以$|G|\le \frac{1}{2}\phi(n)$，所以$|W|\ge\frac{1}{2}\phi(n)$，引理得证。

再考虑Miller-Rabin算法, 设$n$为奇素数，则存在$q,m$使得:
$$
n-1=2^qm(q\le1)
$$
构造序列：
$$
a^m(mod~n), a^{2m}(mod~n), a^{4m}(mod~n),...,a^{2^qm}(mod~n)
$$
由费马定理知，最后一项$a^{n-1}\equiv 1(mod~n)$.而如果$x^2\equiv1(mod~n)$的根是$x=\pm1$称为平凡根，其他根是不平凡的，如果存在非平凡根，那么一定是合数。

所以检测该序列中某项是1，那么前一项必须是$1$或者$n-1$否则$n$是合数。
\begin{algorithm}[h]
    \SetAlgoLined
    \KwIn{$n$}
    \KwOut{$n$是合数或者素数}
    \BlankLine
    $q\leftarrow 0, m \leftarrow n-1$\\
    \Repeat{$m$是奇数}{
        $m\leftarrow \frac{m}{2}$\\
        $q\leftarrow q + 1$\\
    }

    $a\leftarrow $Random(2, $n-1$)\\
    $x_0\leftarrow$EXPMOD($a, m, n$)\tcp{$x_0=a^m(mod~n), O(\log^3n)$}
    \For{$i\leftarrow1$\KwTo $q$}{
        $x_i\leftarrow x_{i-1}^2(mod~n)$\tcp{$O(\log^2n)$}
        \If{$x_i=1$ and $x_{i-1}\neq n-1$ and $x_{i-1}\neq 1$}{
            \Return{合数}\tcp{100\% correct}
        }
    }
    \If{$x_q\neq 1$}{
        \Return{合数}\tcp{100\% correct}
    }
    \Return{素数}\tcp{\textbf{not} 100\% correct}
    \caption{单次的Miller-Rabin素数判定算法}
\end{algorithm}

每次测试出错(把合数判定为素数)概率至多是$\frac{1}{2}$, 这样我们重复运行$k$次，就可以将出错的概率降低到$2^{-k}$, 那么取$k=\lceil \log n\rceil$，失败的概率不超过$2^{-\lceil \log n\rceil}\le \frac{1}{n}$, 而这就是\textbf{Miller-Rabin算法}.


\subsubsection{验证串的相等性}
对于一个比特串$w$，设$I(w)$是$w$表示的一个整数，那么可以选一个素数$p$，这样得到一个指纹函数:
$$
I_p(x)=I(x)(mod~x)
$$
对于$I_p(x)=I_p(y)$，而$x\neq y$的现象，我们称之为\textbf{假匹配}。每次验证时重新选择$p$，并重复选择$p$, 可以增强$x=y$的可信度。

我们知道，设$\pi(n)$是小于$n$的不同素数的个数，那么$\pi(n)$趋近于$\frac{n}{\ln n}$.
对于两个长度为$n$的串$x,y$，如果发生假匹配，那么$p|(I(x)-I(y))$，所以其概率为：
$$
\frac{|\{p\le M|p is prime and p | I(x)-I(y)\}|}{\pi(M)}\le\frac{\pi(n)}{\pi(M)}
$$

\noindent \textbf{关键观察}：假匹配发生时，素数 $p$ 必须满足 $p \mid (I(x) - I(y))$。由于 $x$ 和 $y$ 是长度为 $n$ 的比特串，$I(x)$ 和 $I(y)$ 的绝对值不超过 $2^n - 1$，因此
$$
|I(x) - I(y)| \leq 2^n - 1.
$$
(这个结论需要在$n>11$后成立，这是由切比雪夫函数完成估计的，最终$n$以内的素数的乘积的增长速率约为$e^{0.9n}$，是高于$2^n-1$这个速率的，所以该估计是正确的)

而如果取$M=n^2$，重复执行$k$次算法，这样最后概率最大是$(\frac{2}{n})^k$.

\subsubsection{模式匹配}
确定模式串是否出现在文本串中(假设都是0,1串)。利用随机算法比较模式指纹$I_p(Y)$和文本块$I_p(X(j))$，让文本块的指纹要易于计算，最好还可以通过$X(j)$计算出$X(j+1)$

仿照上一题是思想，我们把这里的指纹就假定为选定的串对应的数的大小，那么有：
$$
I_p(X(j+1))=(2I_p(X(j)) -2^mx_j+x_{j+m})(mod~p)\\
$$

\begin{algorithm}[h]
    \KwIn{长度为$n$的文本串$X$，长度为$m$的模式串$Y$}
    \KwOut{模式串$Y$在文本串$X$中的第一个匹配位置，未匹配返回0}
    \BlankLine

    在小于$M$的素数中随机选择素数$p$\\
    $j\leftarrow1$\\
    计算$I_p(Y), I_p(X_j)$\\
    \While{$j\le n - m + 1$}{
        \If{$I_p(Y)=I_p(X_j)$}{
            \Return $j$
        }
        $I_p(X_{j+1})\leftarrow f(I_p(X_j),p)$\\
        $j\leftarrow j+1$
    }
    \Return 0
    \caption{Rabin-Karp算法}
\end{algorithm}

易知该算法的时间复杂度是$O(m+n)$

\subsubsection{货郎担问题}
假设有一个旅行商人要拜访n个城市，路径选择需满足以下限制：
\begin{enumerate}
    \item 每个城市只能拜访一次；
    \item 最终必须回到出发城市。
\end{enumerate}

​优化目标​：选择总路程最短的路径

\begin{algorithm}[H]
\SetAlgoLined
\KwIn{城市数 $n$，每两城市间距离矩阵 $D$}
\KwOut{最优路径 $bestPath$，路径总长 $bestDistance$}
\BlankLine
$bestPath \leftarrow [\,]$ \tcp{初始化最优路径为空}
$bestDistance \leftarrow \infty$\tcp{初始化最短距离为无穷大}
$path \leftarrow \text{随机排列}(1, 2, \dots, n)$ \tcp{生成初始随机路径}
$distance \leftarrow \text{计算路径总长度}(path, D)$ \tcp{计算初始路径长度}
\BlankLine
\For{$j \leftarrow 1$ \KwTo $iteration$}{
  \If{$distance < bestDistance$}{
    $bestPath \leftarrow path$ \tcp{更新最优路径}
    $bestDistance \leftarrow distance$ \tcp{更新最短距离}
  }
  $path \leftarrow \text{生成新随机路径}()$ \tcp{重新采样路径}
  $distance \leftarrow \text{计算路径总长度}(path, D)$ \tcp{重新计算长度}
}
\Return{$bestPath$, $bestDistance$} \tcp{返回找到的最优解}
\caption{随机采样求解TSP算法 (RandomTSP)}
\end{algorithm}

\subsection{随机游走算法}
常用有限马尔科夫链来分析随机游走算法。

马尔可夫特性：随机过程在给定t时刻的状态为Xt，该过程的后续状态及其出现的概率与t之前的状态无关。也就是说，过程当前的状态包括了过程所有的历史信息，该过程的进一步发展完全由当前状态所决定，与当前状态之前的历史无关，这种性质也称为\textbf{无后效性}或\textbf{无记忆性}。

马尔科夫不等式：设$X$为只取非负的随机变量，数学期望$E(X)=\mu$, 对于任意$\varepsilon>0, P(X\geq\varepsilon)\le\frac{\mu}{\varepsilon}$

\begin{enumerate}
    \item \textbf{Integral method} (for continuous case):
    $$
    E(X) = \int_0^\infty x f(x)dx \geq \int_\varepsilon^\infty x f(x)dx \geq \int_\varepsilon^\infty \varepsilon f(x)dx \ge \varepsilon P(X \geq \varepsilon)
    $$
    
    \item \textbf{Indicator method} (general):
    Define $Y = \varepsilon \cdot \mathbf{1}_{\{X \geq \varepsilon\}}$, then:
    $$
    E(Y) = \varepsilon P(X \geq \varepsilon) \leq E(X)
    $$
\end{enumerate}

设一个离散的随机序列$\{X_0, X_1, ...\}$，其中每个$X_i$都从有限集中取值，若它满足
$$
P(X_{i+1}=x_{i+1}|X_i=x_i,X_{i-1}=x_{i-1},...,X_0=x_0)=P(X_{i+1}=x_{i+1}|X_i=x_i)=p_{x_i, x_{i+1}}
$$

这称之为\textbf{有限马尔科夫链}

对于有限马氏链，不妨设$X_t$取值的状态空间为$\{0,1,\ldots,n\}$，于是$p_{i,j}$（从状态$i$转移到状态$j$的概率）可组成一个$n+1$阶方阵
$$
P = [p_{i,j}],
$$
称为\textbf{一步转移矩阵}，每一行元素之和等于$1$。

设$p_i(t)$表示在时刻$t$处于状态$i$的概率，则
$$
p_i(t) = p_0(t-1)p_{0i} + p_1(t-1)p_{1i} + \cdots + p_{n-1}(t-1)p_{(n-1)i} + p_n(t-1)p_{ni}
$$

设$\boldsymbol{p}(t)$表示向量$(p_0(t), p_1(t), \ldots, p_n(t))$，有
$$
\boldsymbol{p}(t) = \boldsymbol{p}(t-1)P
$$

对于任意$m$，定义$m$步转移矩阵
$$
P^{(m)} = \big[p_{i,j}^{(m)}\big],
$$
其中
$$
p_{i,j}^{(m)} = P(X_{t + m} = j \mid X_t = i).
$$
因此，
$$
P^{(m)} = P^m,\\
\boldsymbol{p}(t + m) = \boldsymbol{p}(t) P^{m}.
$$

\subsection{二元布尔可满足问题}

\begin{algorithm}[H]
\SetAlgoLined
\KwIn{整数$n$（变元数量），整数$m$（子句数量），合取范式$F$（每个子句至多两个文字）}
\KwOut{可满足赋值，或宣布不可满足}
\BlankLine
随机初始化赋值$\sigma$：对每个变元$x_i$随机赋$\text{True}$或$\text{False}$\;
\For{$k \leftarrow 1$ \textbf{to} $2mn^2$}{
    \If{$\sigma$满足$F$}{
        \Return{$\sigma$} \tcp*{找到可满足解}
    }
    随机选择$F$中一个不满足的子句$C$\;
    从$C$中随机选择一个文字$l$\;
    翻转$l$对应变元的赋值 \tcp*{将$\text{True}$改为$\text{False}$或反之}
}
\Return{"不可满足"} \tcp*{超过最大迭代次数未找到解}
\caption{2-SAT随机局部搜索算法}
\end{algorithm}

\noindent\textbf{定理：} 若输入公式是不可满足的，则上述随机游走算法正确宣布不可满足；若输入公式是可满足的，则上述算法以 $1-\frac{1}{2^m}$ 概率找到可满足赋值。

定义马尔可夫链$\{Y_0,Y_1,Y_2,\cdots\}$如下：
\begin{align*}
Y_0 &= X_0 \\
P(Y_{t+1}=1|Y_t=0) &= 1 \\
P(Y_{t+1}=j+1|Y_t=j) &= \frac{1}{2} \\
P(Y_{t+1}=j-1|Y_t=j) &= \frac{1}{2} \quad (j > 0)
\end{align*}

显然有：
\begin{align*}
P(Y_{t+1}=j+1|Y_t=j) &= \frac{1}{2} \geq P(X_{t+1}=j+1|X_t=j) \\
P(Y_{t+1}=j-1|Y_t=j) &= \frac{1}{2} \geq P(X_{t+1}=j-1|X_t=j)
\end{align*}

从任意状态出发，$Y_t$比$X_t$花费更长的期望时间进入状态$n$。因此$Y_t$进入状态$n$的时间作为算法期望运行时间的上界。

设$h_j$表示从状态$j$到状态$n$的期望运行时间，则有：
\begin{align*}
h_n &= 0 \\
h_0 &= h_1 + 1 \\
h_j &= \frac{h_{j+1}+1}{2} + \frac{h_{j-1}+1}{2}, \quad 0 < j < n
\end{align*}

由归纳法可证，对于所有$j$，有：
$$
h_j \leq n^2 - j^2 \leq n^2
$$

因此有：
$$
E[\text{$X$从$X_0$到达$n$的时间}] \leq E[\text{$Y$从$Y_0$到达$n$的时间}] \leq n^2
$$

取$\varepsilon = 2n^2$，根据马尔可夫不等式，则有：
$$
P(\text{$X$从$X_0$到达$n$的时间} \geq 2n^2) \leq \frac{n^2}{2n^2} = \frac{1}{2}
$$

\noindent\textbf{证明：}
\begin{itemize}
    \item[$\bullet$] \textbf{不可满足情形：} 显然，算法无法找到可满足赋值，因此正确宣布不可满足。
    
    \item[$\bullet$] \textbf{可满足情形：} 算法有可能找不到可满足赋值而出错，下面证出错概率为 $\frac{1}{2^m}$。
\end{itemize}

    
\begin{enumerate}
    \item \textbf{汉明距离定义}：设$\sigma^*$为最优赋值，当前赋值$\sigma$与$\sigma^*$的汉明距离$d(\sigma,\sigma^*)$表示两者取值不同的变量个数。
    
    \item \textbf{改进概率}：对每个不满足子句，至少以$\frac{1}{2}$概率减少$d(\sigma,\sigma^*)$（因子句最多2个文字）。
    
    \item \textbf{收敛分析}：经过$2n^2m$步迭代，未到达$\sigma^*$的概率：
    $$
    P(\text{失败}) \leq \left(1-\frac{1}{2n}\right)^{2n^2m} \approx e^{-nm} \leq \frac{1}{2^m}
    $$
\end{enumerate}



\subsection{本章作业}
\subsubsection{优惠劵收集}
有$n$种优惠券，每种优惠券数目不限。每次实验随机抽取一张优惠券，优惠券来自$n$种中任何一种的概率均等。计算需要多少次实验才能使得收集的优惠券包含$n$种中每一种至少一张。

设$X_i$表示在已收集$i-1$种优惠券后，首次抽到第$i$种新优惠券所需的次数（$i=1,2,\dots,n$）。 

在已收集$k$种优惠券时，抽到新优惠券的概率为$p_k = \frac{n-k}{n}$，因此$X_{k+1}$服从参数为$p_k$的几何分布，其期望为：  
$$
   E(X_{k+1}) = \frac{1}{p_k} = \frac{n}{n-k}
$$

集齐所有$n$种优惠券的总次数$T$为各阶段次数之和：  
   $$
   T = X_1 + X_2 + \dots + X_n
   $$
   总期望为：  
   $$
   E(T) = \sum_{k=0}^{n-1} \frac{n}{n-k} = n \sum_{k=1}^n \frac{1}{k}=nH_n
   $$
当$n$较大时，$H_n \approx \ln n + \gamma$。因此：  
   $$
   E(T) \approx n \ln n + \gamma n
   $$
集齐$n$种优惠券的平均实验次数为：  
$$
E(T) = n H_n = n \left(1 + \frac{1}{2} + \dots + \frac{1}{n}\right)
$$
对于大$n$，可近似为$n \ln n + \gamma n, \text{其中}\gamma\approx 0.577$，

\subsubsection{用公平硬币随机序列生成}
假定你有一枚公平硬币，请设计一个有效的随机算法用来生成整数$1,2,...,n$的一个随机排列。并分析你的算法的时间复杂性.


\begin{algorithm}[H]
\SetAlgoLined

\KwIn{整数$n$（排列长度）}
\KwOut{数组$A[1..n]$包含$1$到$n$的随机排列}
\BlankLine
初始化数组$A$：$A[i] \leftarrow i$ 对 $i=1$到$n$\;
\For{$i \leftarrow n$ \textbf{downto} $2$}{
  生成$\lceil \log_2 i \rceil$次公平硬币投掷，其中正面朝上为1，反面为0来得到随机数的二进制串，得到随机数$j \in [0, 2^{\lceil \log_2 i \rceil}-1]$\;
  \While{$j \geq i$}{
    重新生成$j$ \tcp*{拒绝采样保证均匀性}
  }
  交换$A[i]$与$A[j+1]$\;
}
\Return{$A$}
\caption{Fisher-Yates随机排列算法（基于公平硬币）}
\end{algorithm}

考虑$i$次枚举后才找到
$$
\sum_{i=1}^\infty q^i=\frac{q}{1-q}=1-p\\
解得q=1-\frac{1}{2-p}\\
$$


在最优情况下，所有的$j$都被接受，最优时间复杂度是
\begin{equation*}
    \sum_{i=2}^n\lceil\log_2i\rceil\triangleq C_{opt}(n)\\
\end{equation*}

有：

$$
C_{opt}(n)\le \sum_{i=2}^n\log_2n=O(n\log n)
$$

也有：(此处省略分部积分的步骤)
$$
C_{opt}(n)\ge\sum_{i=2}^n\log_2 i=\sum_{i=2}^n\int_{i-1}^{i}\log_2idx\ge\sum_{i=2}^n\int_{i-1}^{i}\log_2xdx=\int_{1}^n\log_2xdx=O(n\log n)
$$
因此由夹逼定理有：$C_{opt}(n)=O(n\log n)$

下面考虑最坏情况$C_{wor}(n)$


对于$i$情况下单次$j$的取值，拒绝概率为：

\begin{equation*}
    P_{reject}(i)=\frac{2^{\lceil\log_2i\rceil}-i}{2^{\lceil\log_2i\rceil}}=1-\frac{i}{2^{\lceil\log_2i\rceil}}\triangleq p_i
\end{equation*}


考虑$i$情况下$k$次$j$的取值被拒绝的概率是$p_i^k$，每次被拒绝都会多进行$\lceil\log_2i\rceil$次硬币投掷实验，因此$k$次取值后产生的次数为：$p_{i}^{k}\cdot k\lceil\log_2i\rceil$，因此$i$情况下的总投掷次数为：（跳过等差乘等比级数的数学推导）
$$
w_i=\sum_{k=1}^\infty p_{i}^{k}\cdot k\lceil\log_2i\rceil=\frac{p_i(1+p_i)}{(1-p_i)^2}\lceil\log_2i\rceil=\frac{(2^{\lceil\log_2i\rceil}-i)(2^{\lceil\log_2i\rceil+1}-i)}{i^2}\le\frac{i\cdot3i}{i^2}=3
$$
因为拒绝而多出的总次数为：
$$
C_{penalty}(n)=\sum_{i=2}^nw_i\le 3n=O(n)
$$
所以最坏时间复杂度为：
$$
C_{wor}(n)=C_{penalty}(n)+C_{opt}(n)=O(n\log n)\\
$$
综上，该算法的时间复杂度是$O(n\log n)$
(当然可以改成如果$j>i$，就$j \%=i$，这样也就省了这步计算了)


\subsubsection{矩阵互逆}
给定两个$n \times n$方阵$A$和$B$，判定它们是否互逆（即$AB = I$）。算法基于蒙特卡罗方法，通过随机采样验证$AB$矩阵元素的性质：
\begin{itemize}
    \item 对角线元素应为$1$（允许误差$\epsilon$）
    \item 非对角线元素应为$0$（允许误差$\epsilon$）
\end{itemize}

\begin{algorithm}[H]
\SetAlgoLined
\KwIn{矩阵 $A$, $B$，误差阈值 $\epsilon$，采样次数 $T$}
\KwOut{YES 或 NO}
\For{$t \leftarrow 1$ \KwTo $T$}{
    $J \leftarrow \text{rand}(1, n)$ \tcp{随机选行}
    $K \leftarrow \text{rand}(1, n)$ \tcp{随机选列}
    $\text{ans} \leftarrow 0$
    \For{$l \leftarrow 1$ \KwTo $n$}{
        $\text{ans} \leftarrow \text{ans} + A[J][l] \times B[l][K]$\\
    }
    \If{($J = K$ 且 $|\text{ans} - 1| > \epsilon$) 或 ($J \neq K$ 且 $|\text{ans}| > \epsilon$)}{
        \Return{NO}
    }
}
\Return{YES}
\end{algorithm}

\begin{itemize}
    \item \textbf{时间复杂度}：$O(T \cdot n)$
    \begin{itemize}
        \item 每次采样需要$O(n)$时间计算矩阵元素
        \item 共进行$T$次采样
    \end{itemize}
    \item \textbf{空间复杂度}：$O(1)$（仅需存储临时变量）
\end{itemize}

\subsubsection{三元可满足问题}
请设计随机游走算法求解三元可满足问题，并分析算法的期望运行时间。

\begin{algorithm}[H]
\SetAlgoLined
\KwIn{3-SAT公式 $F$，最大迭代次数 $T$}
\KwOut{满足 $F$ 的赋值（若找到），否则返回“无解”}
随机初始化变量赋值 $\mathbf{x} = (x_1, \dots, x_n)$\;
\For{$t \leftarrow 1$ \KwTo $T$}{
    \If{$\mathbf{x}$ 满足 $F$}{
        \Return{$\mathbf{x}$}\;
    }
    随机选择一个不满足的子句 $C$\;
    随机选择 $C$ 中的一个变量 $x_j$\;
    翻转 $x_j$ 的值（$\text{True} \leftrightarrow \text{False}$）\;
}
\Return{“无解”}\;
\end{algorithm}

\begin{itemize}
    \item \textbf{递推关系}（3文字子句）：
    $$
    h_j = \frac{2}{3}h_{j+1} + \frac{1}{3}h_{j-1} + 1
    $$
    
    \item \textbf{特征根}：
    $$
    r_1 = 1, \ r_2 = -\frac{1}{2} \implies h_j \approx 3^j
    $$
    
    \item \textbf{指数下界}：
    $$
    E[T] \geq 3^n - 1 = \Omega(3^n)
    $$
\end{itemize}


\subsubsection{拉斯维加斯型与蒙特卡洛型的期望运行时间推导关系}

设$A$是一个蒙特卡洛随机算法，满足：
\begin{itemize}
    \item 期望运行时间不超过$T(n)$
    \item 返回正确解的概率为$p(n)$
    \item 任何解的正确性可在$T'(n)$时间内验证
\end{itemize}

我们构造拉斯维加斯算法$A'$如下：
\begin{enumerate}
    \item 运行$A$（耗时$\leq T(n)$）
    \item 验证解的正确性（耗时$\leq T'(n)$）
    \item 若正确则返回解, 否则重复验证
\end{enumerate}

$$
   E[\text{时间}] \leq \frac{T(n) + T'(n)}{p(n)}
$$

期望时间公式成立因为：
  $$
  E[\text{总时间}] = \sum_{k=1}^{\infty} \left[(1-p(n))^{k-1} p(n) \cdot k (T(n) + T'(n))\right] = \frac{T(n) + T'(n)}{p(n)}
  $$
该转换在保证结果正确性的前提下，仅增加了多项式时间开销
\end{document}